\documentclass[conference]{IEEEtran}

% ===================== Packages =====================
\usepackage{amsmath,amssymb}
\usepackage{graphicx}
\usepackage{booktabs}
\usepackage{multirow}
\usepackage{adjustbox}
\usepackage[hidelinks]{hyperref}

% ===================================================
\begin{document}

% ===================== TITLE =====================
\title{A Comparative Study of Gradient and Non-Gradient Optimisation Methods for Fuzzy Inference Systems}

% ===================== AUTHORS  =====================
\author{}
\maketitle

% ===================== ABSTRACT =====================
\begin{abstract}
Optimisation plays a critical role in the performance of fuzzy inference systems,
yet empirical comparisons between optimisation methods are often confounded by
constraint handling, membership-function design, and inconsistent configuration
selection protocols.
This paper presents a staged and reproducible experimental framework for
Mamdani-type fuzzy inference systems (FIS) within a unified and fully
differentiable PyTorch implementation, with the practical goal of producing
actionable guidance on optimisation method selection.

The proposed framework explicitly separates early design decisions from
optimiser-specific configuration search and final benchmarking.

In Stage~I, we provide principled design guidance for constraint parameterisation
and membership-function granularity initialisation.
Based on established theoretical insights from prior work, a CASPs variant is
selected according to optimisation objectives, and a lightweight structural evaluation is performed to identify a representative
membership-function granularity for each dataset without performing full optimisation.
This stage serves to reduce the design space and avoid confounding downstream
results with arbitrary structural choices.

In Stage~II, with constraint handling and membership-function granularity fixed,
optimiser-specific configuration selection is performed over candidate
initialisation strategies and hyperparameter settings using a unified successive
halving (SH) protocol under limited optimisation budgets.
This ensures fair and efficient allocation of computational resources across
both gradient-based and non-gradient optimisation methods.

In Stage~III, the resulting optimiser--configuration combinations are evaluated
on held-out test sets across ten benchmark regression datasets using multiple
random seeds.
All datasets, model architectures, and evaluation budgets are held constant
across optimisers, enabling controlled comparison in terms of predictive
accuracy, robustness, and computational efficiency.

By explicitly decoupling constraint selection, structural initialisation, and
configuration optimisation, the proposed framework reduces methodological
confounding and provides reproducible, system-level guidance for optimising
differentiable Mamdani-type fuzzy inference systems.
\end{abstract}



\begin{IEEEkeywords}
Fuzzy inference systems, optimisation, gradient-based optimisation,
non-gradient optimisation, evolutionary algorithms, empirical comparison
\end{IEEEkeywords}

% ===================================================
\section{Introduction}

Since the introduction of fuzzy sets by Zadeh in 1965~\cite{zadeh1965},
fuzzy inference systems (FIS) have become a widely used modelling paradigm
for regression, control, and decision-support, particularly in applications
where interpretability and linguistic transparency are required.
A practical FIS relies on tunable components such as membership functions
and rule parameters, and its predictive performance depends critically on
how effectively these parameters are optimised.

Historically, fuzzy-system optimisation has been dominated by non-gradient
and evolutionary methods, motivated by highly non-convex search landscapes
and the presence of non-differentiable design choices.
Representative approaches include genetic algorithms (GA)~\cite{holland1975ga},
particle swarm optimisation (PSO)~\cite{kennedy1995pso}, and differential
evolution (DE)~\cite{storn1997de}.
More recently, advances in differentiable programming have enabled fully
differentiable implementations of Mamdani-type fuzzy inference systems,
making gradient-based optimisation increasingly practical.
In particular, adaptive first-order methods such as Adam~\cite{kingma2015adam}
and RMSprop~\cite{tieleman2012rmsprop} have become standard tools for
large-scale continuous optimisation.

Despite the widespread use of both optimisation paradigms, there remains
limited consensus on their relative effectiveness under controlled and
consistent experimental conditions for Mamdani-type fuzzy inference systems.
Many existing studies introduce new optimisation variants or hybrid strategies
and evaluate them under heterogeneous datasets, architectures,
constraint-handling mechanisms, membership-function designs, and tuning
protocols.
As a result, observed performance differences are difficult to attribute
unambiguously to the optimisation method itself, and conclusions may be
sensitive to arbitrary design choices made early in the modelling pipeline.

Recent work has demonstrated that gradient-based optimisation of Mamdani-type
fuzzy systems can be implemented efficiently via automatic differentiation
within modern machine learning frameworks~\cite{chen2024gradient}, while
non-gradient approaches remain competitive in settings characterised by
highly non-convex or multimodal error surfaces~\cite{zhou2023stochastic}.
In parallel, constraint-aware parameterisations such as Constraints Always
Satisfied Parameters (CASPs) provide a principled mechanism for enforcing
semantic and geometric constraints during optimisation without resorting to
ad-hoc projection or penalty schemes~\cite{chen2025casp}.

However, in many comparative studies, constraint handling, membership-function
granularity, parameter initialisation, and optimiser hyperparameters are
treated as a single entangled design space.
This entanglement makes it difficult to isolate the effect of the optimisation
method itself and increases the risk of methodological confounding.
In particular, early structural design choices—such as the degree of
constraint strictness and the number of membership functions per variable—can
strongly influence downstream optimisation behaviour, yet are often selected
heuristically or fixed without explicit justification.

To address these limitations, this paper proposes a controlled and staged
experimental framework for the empirical comparison of gradient-based and
non-gradient optimisation methods for Mamdani-type fuzzy inference systems
within a unified and fully differentiable implementation.
Rather than jointly optimising all design choices, the framework explicitly
separates early design guidance, optimiser-specific configuration selection,
and final performance evaluation.

Specifically, the study adopts a three-stage experimental protocol that
explicitly separates early design guidance, optimiser-specific configuration
selection, and final performance evaluation.
This separation is intended to minimise methodological confounding and to
enable interpretable comparison of optimisation behaviour under controlled
conditions.



\section{Related Work}

Optimisation has long been a central research topic in the development of
fuzzy inference systems.
Early work predominantly focused on non-gradient and evolutionary algorithms,
motivated by highly non-convex search landscapes and the frequent presence of
non-differentiable operators.
Genetic algorithms (GA), particle swarm optimisation (PSO), differential
evolution (DE), and covariance matrix adaptation evolution strategy (CMA-ES)
have been widely applied to optimise membership functions, rule parameters,
and fuzzy partitions
\cite{holland1975ga, kennedy1995pso, storn1997de, hansen2001cmaes}.
Comprehensive surveys further document the effectiveness of evolutionary
approaches across a broad range of fuzzy-system applications
\cite{cordon2001review}.
However, empirical comparisons in this literature are often conducted under
heterogeneous architectures, initialisation choices, and tuning protocols,
making it difficult to isolate the effect of the optimisation method itself.

More recently, advances in differentiable programming have enabled
gradient-based optimisation of Mamdani-type fuzzy inference systems.
By expressing fuzzy inference operations as differentiable tensor computations,
automatic differentiation can be used to optimise membership-function
parameters efficiently within modern machine learning frameworks
\cite{chen2024gradient}.
While gradient-based approaches have demonstrated competitive performance,
their behaviour remains sensitive to constraint handling, parameter
initialisation, and structural design choices, which are not always treated
explicitly in comparative studies.

Constraint handling constitutes a particularly important aspect of fuzzy
system optimisation, as membership-function parameters are subject to semantic
and geometric constraints that must be preserved throughout optimisation.
Traditional approaches often enforce such constraints through heuristic
projection, penalty terms, or restricted parameterisations.
Constraint-aware reparameterisation methods, such as Constraints Always
Satisfied Parameters (CASPs), provide a principled alternative by guaranteeing
constraint satisfaction by construction
\cite{chen2025casp}.
The CASPs framework introduces multiple variants that trade off constraint
strictness, interpretability, and optimisation flexibility.
Rather than re-evaluating these trade-offs experimentally, the present study
adopts CASPs as a constraint-handling mechanism and follows the design
guidelines established in prior work.

Beyond constraint handling, the structural design of fuzzy inference systems,
including the number of membership functions per input and output variable,
plays a critical role in determining model capacity and optimisation behaviour.
Nevertheless, membership-function granularity is often fixed a priori or
implicitly tuned together with optimiser hyperparameters, making its influence
difficult to disentangle from optimiser performance.
This structural aspect is rarely addressed explicitly in large-scale
comparative optimisation studies.

A further methodological challenge in comparative optimisation research is
hyperparameter tuning fairness.
Inconsistent or optimiser-specific tuning practices can substantially bias
empirical results and lead to misleading conclusions.
Multi-fidelity approaches such as successive halving provide a principled
framework for allocating computational budget across competing configurations
\cite{jamieson2016sh}, yet remain underutilised in the context of fuzzy inference
system optimisation.

Motivated by these gaps, this paper revisits the empirical comparison of
gradient-based and non-gradient optimisation methods for Mamdani-type fuzzy
inference systems under a unified and fully differentiable implementation.
By explicitly separating constraint parameterisation, structural
initialisation, configuration selection, and final evaluation, the proposed
framework aims to reduce methodological confounding and provide reproducible,
system-level guidance on optimiser behaviour.

% ===================================================
\section{Methodology}\label{sec:method}

This section describes the fuzzy inference system model, optimisation objective,
and the staged experimental framework adopted in this study.
The methodology is explicitly designed to separate early design decisions from
optimiser-specific configuration search and final performance evaluation, thereby
reducing methodological confounding.
All stages described in this section correspond directly to the implemented
experimental pipeline.

% ===================================================
\subsection{Fuzzy Inference System Model}

We consider a Mamdani-type fuzzy inference system (FIS) with $d$ input variables
and a single output variable.
For each input variable $x_k$, a set of trapezoidal membership functions (MFs)
is defined, and fuzzy rules are constructed by combining antecedent MFs across
all input dimensions.

Let $\theta$ denote the concatenation of all antecedent membership-function
parameters.
In this study, optimisation is restricted exclusively to antecedent MF
parameters.
The rule structure and consequent output MFs are constructed in an initial
rule-generation stage and remain fixed throughout all optimisation experiments.
This design choice isolates optimiser behaviour on continuous antecedent
parameter learning and avoids confounding effects introduced by joint
optimisation of consequents or rule structures.

\paragraph*{Rule base construction and MF initialisation.}
For each dataset and random seed, the fuzzy rule base is constructed once
prior to optimisation and remains fixed throughout all experimental stages.
A k-means-based discretisation procedure is applied independently to each
input variable and to the output variable, producing ordered categorical
labels based on cluster means.
Fuzzy rules are generated by enumerating observed combinations of input
categories, and each rule is assigned a consequent corresponding to the most
frequent output category within the associated data subset.

Initial trapezoidal membership-function parameters are derived directly from
the empirical ranges of the resulting clusters using a deterministic
range-based construction.
This procedure is identical across all optimisation methods and random seeds,
ensuring that observed performance differences arise solely from optimiser
behaviour rather than differences in rule structure or initialisation.

All fuzzy inference operations, including MF evaluation, rule firing,
aggregation, and defuzzification, are implemented as differentiable tensor
operations within a unified PyTorch framework.
This enables the use of both gradient-based and non-gradient optimisation
methods under an identical forward-pass and loss-evaluation pipeline.
% ===================================================
\subsection{Optimisation Objective}

Given a training dataset $\{(x_i, y_i)\}_{i=1}^{N}$, optimisation is formulated
as minimisation of the mean squared error (MSE):
\begin{equation}
\mathcal{L}(\theta)
=
\frac{1}{N}
\sum_{i=1}^{N}
\left( \hat{y}(x_i) - y_i \right)^2 ,
\end{equation}
where $\hat{y}(x_i)$ denotes the crisp output produced by centroid
defuzzification.
All optimisation methods considered in this study minimise the same objective
over an identical parameter space.

% ===================================================
\subsection{Stage I: Constraint and Initialisation Design Guidance}

Stage~I establishes principled design guidance for
(i) constraint parameterisation and
(ii) membership-function granularity initialisation.
Stage~I reduces the design space prior to optimiser-specific
configuration search, ensuring that subsequent comparisons reflect optimiser
behaviour rather than arbitrary structural choices.

% ---------------------------------------------------
\subsubsection{Constraint Parameterisation via CASPs}

Membership-function parameters in fuzzy inference systems are subject to semantic
and geometric constraints, such as ordering and non-negativity, which must be
preserved throughout optimisation.
This study adopts the Constraints Always Satisfied Parameters (CASPs) framework
\cite{chen2025casp}, which enforces such constraints by construction through
reparameterisation.

The CASPs framework defines multiple variants that differ in the degree of
constraint coupling across membership functions.
CASPs-Single imposes the strictest constraints by tightly coupling adjacent MFs,
ensuring well-structured overlap and strong semantic interpretability.
However, this strict coupling substantially restricts parameter freedom and can
limit optimisation performance, particularly in complex or highly non-convex
loss landscapes.

CASPs-Adapted relaxes some of the inter-MF coupling constraints while preserving
basic structural regularity.
This variant provides a compromise between interpretability and optimisation
flexibility, but may still introduce structural bias that affects optimiser
behaviour.

CASPs-Free enforces only the minimal internal ordering constraints required for
each MF to remain valid.
By removing inter-MF overlap constraints, CASPs-Free provides maximal parameter
freedom during optimisation.
Prior work has shown that this variant consistently achieves superior RMSE
performance when predictive accuracy is prioritised over strict semantic
structure \cite{chen2025casp}.

In this study, the primary objective is to compare optimisation methods in terms
of predictive performance under controlled and reproducible conditions.
Accordingly, CASPs-Free is adopted uniformly throughout all subsequent stages.
Nevertheless, CASPs-Single or CASPs-Adapted may be preferable in applications
where interpretability, smooth linguistic transitions, or strict semantic
structure are prioritised over raw predictive accuracy.

% ---------------------------------------------------
\subsubsection{Membership-Function Granularity Initialisation}

Beyond constraint parameterisation, the number of membership functions per input
and output variable constitutes a key structural design choice that directly
affects model capacity and optimisation behaviour.
Overly coarse MF partitions may limit approximation capability, whereas
excessively fine partitions can lead to redundant rules, dominant rule coverage,
and unstable or inefficient optimisation.

To determine appropriate MF granularities without performing parameter
optimisation, a lightweight, structure-only evaluation is conducted in
Stage~I.
Candidate MF-count configurations are randomly sampled within predefined ranges
for each dataset.
For each candidate configuration, a fuzzy rule base is constructed using the
same k-means-based discretisation procedure adopted throughout the study, while
all membership-function parameters remain fixed.

Each configuration is assessed using a composite structural score that captures
the suitability of the induced rule base, rather than predictive performance.
The score incorporates:
(i) a lower bound on achievable error derived from within-rule output variance,
(ii) rule coverage and compactness statistics,
(iii) entropy-based measures of rule and membership-function utilisation, and
(iv) penalties for dominant or redundant rules.

Crucially, no parameter optimisation or iterative training is performed during
this stage.
The purpose of this evaluation is solely to exclude structurally degenerate or
unnecessarily complex MF configurations and to identify reasonable,
dataset-specific MF granularities.
The selected MF counts are then fixed and recorded in the experimental
configuration registry.
All subsequent experiments in Stage~II and Stage~III are conducted using these
fixed MF granularities without further modification.
\paragraph*{MF granularity selection protocol.}
Concretely, MF granularity selection is implemented as an optimiser-independent
structural evaluation procedure.
For each dataset, candidate MF-count configurations are randomly sampled from
bounded discrete ranges for all input and output variables.
Each sampled configuration is evaluated by constructing a rule base via
one-dimensional k-means discretisation and computing a set of fixed
rule-based structural metrics, including a lower bound on achievable RMSE,
rule coverage statistics, entropy-based utilisation measures, and penalties
for dominant or redundant rules.

Configurations violating predefined structural constraints are discarded
without further consideration.
The remaining configurations are ranked using a fixed weighted scoring
function defined over these structural metrics.
Based on this ranking, a single MF configuration is selected for each dataset
and fixed for all subsequent experimental stages.
No optimiser-specific information is used during this process, and no parameter
optimisation is performed at any point in Stage~I.


% ===================================================
\subsection{Stage II: Configuration Selection under Fixed Design}

In Stage~II, optimiser-specific configuration selection is performed under the
fixed design choices established in Stage~I.
Constraint handling is fixed to CASPs-Free, and MF granularity is restricted to
the ranges identified by the structural evaluation.

For each optimiser, a discrete set of candidate initialisation strategies and
hyperparameter configurations is defined.
Configuration selection is performed using a unified successive halving (SH)
protocol \cite{jamieson2016sh}, which allocates computational budget efficiently
across competing configurations.

All candidate configurations are initially evaluated under limited optimisation
budgets.
Poorly performing configurations are progressively discarded, and the remaining
configurations are re-evaluated under increased budgets.
Throughout Stage~II, performance measurements are used exclusively for
configuration selection and do not contribute to final performance reporting.

% ===================================================
\subsection{Stage III: Final Evaluation Protocol}

Final evaluation is conducted using the optimiser-specific configurations
selected in Stage~II.
Experiments are performed on a subset of ten benchmark regression datasets drawn
from the collection commonly adopted in prior comparative studies
\cite{zhou2023stochastic}.

For each dataset, all optimisers are evaluated under identical experimental
conditions using ten fixed random seeds.
All datasets are normalised using min--max scaling and split into training and
test sets.
The test set is accessed exclusively at this stage and only once per
configuration.

Performance is assessed in terms of predictive accuracy, robustness across
datasets, and computational efficiency.
By restricting Stage~III to final evaluation only, the staged protocol ensures
that reported results reflect optimiser behaviour rather than artefacts of
design selection or configuration tuning.

% ===================================================
% ===================================================
\section{Experimental Evaluation}

Following the staged methodology described in
Section~\ref{sec:method}, this section reports the datasets,
experimental setup, and empirical results obtained under the fixed
design and evaluation protocol.
Stage~I establishes dataset-specific structural design choices,
Stage~II selects optimiser configurations, and Stage~III reports the
final test-set performance used for comparative analysis.


% ===================================================
\subsection{Datasets}

Experiments are conducted on ten benchmark regression datasets selected from
the collection commonly adopted in prior comparative studies such as
\cite{zhou2023stochastic}.
The datasets span a range of sample sizes and input dimensionalities, providing
a diverse testbed for evaluating optimiser behaviour under controlled
conditions.
Table~\ref{tab:datasets} summarises the datasets used in the experimental
evaluation.

\begin{table}[t]
\centering
\caption{Regression datasets used in the experimental evaluation.}
\label{tab:datasets}
\begin{adjustbox}{max width=\linewidth}
\begin{tabular}{lcc}
\toprule
\textbf{Dataset} & \textbf{Samples (N)} & \textbf{Features (d)} \\
\midrule
Abalone   & 4177 & 10 \\
Airfoil  & 1503 & 5  \\
Ankara   & 321  & 9  \\
AutoMPG6 & 84   & 7  \\
Baseball & 1340 & 22 \\
Concrete & 1030 & 8  \\
Izmir    & 1461 & 9  \\
Laser    & 993  & 4  \\
Treasury & 1049 & 15 \\
Wine     & 1599 & 11 \\
\bottomrule
\end{tabular}
\end{adjustbox}
\end{table}

% ===================================================
\subsection{Experimental Setup and Reproducibility}

All experiments are implemented within a unified PyTorch-based framework,
ensuring identical forward-pass computation and loss evaluation across all
optimisers.
For each dataset, data are split into training and test sets using a fixed
five-fold partition, and ten random seeds are used to ensure robustness.
Input variables are normalised using min--max scaling where appropriate.

For each dataset and seed, the fuzzy rule base and initial membership-function
parameters are generated once using a k-means-based discretisation procedure.
Only antecedent membership-function parameters are optimised.
The rule base, MF granularity, constraint parameterisation, and optimisation
budget are held constant across all optimisation methods, ensuring that
observed performance differences can be attributed to optimiser behaviour
rather than confounding design choices.

% ===================================================
\subsection{Stage I: Constraint and Structural Design Decisions}

Stage~I instantiates the design choices established in
Section~\ref{sec:method} for each dataset prior to optimiser comparison.
All decisions made at this stage are fixed for the remainder of the experimental
pipeline to ensure fair and controlled evaluation.

\paragraph*{Constraint parameterisation.}
Based on the design guidelines established in prior work on CASPs
\cite{chen2025casp}, the choice of constraint mode is determined by the primary
optimisation objective.
Since this study focuses on predictive performance rather than strict semantic
regularity, the CASPs-Free parameterisation is adopted uniformly for all
datasets and all subsequent experimental stages.
This ensures maximal parameter flexibility while maintaining validity of
membership functions throughout optimisation.

\paragraph*{Membership-function granularity selection.}
In addition to constraint handling, Stage~I determines the number of membership
functions per input and output variable for each dataset.
Rather than fixing MF counts heuristically, a lightweight structural search is
performed prior to optimisation to exclude degenerate or excessively complex
configurations.

Candidate MF-count configurations are evaluated using rule-based structural
criteria, including coverage statistics, entropy-based utilisation measures,
and a lower bound on achievable error derived from within-rule output variance.
The selected MF granularities are fixed for each dataset and recorded in the
experimental configuration registry.
All optimisation experiments in Stage~II and Stage~III are conducted using these
fixed MF configurations without further modification.

For datasets with high input dimensionality, principal component analysis (PCA)
is applied as a preprocessing step to reduce the input space to five dimensions
prior to membership-function construction, controlling rule-base complexity
while preserving the majority of input variance.


% ===================================================
\subsection{Stage II: Configuration Selection via Successive Halving}

In Stage~II, optimiser-specific configuration selection is performed under the
fixed design choices established in Stage~I.
Constraint handling is fixed to CASPs-Free, and MF granularity is fixed to the
values selected by the structural evaluation.

For each optimiser, a discrete hyperparameter search space is defined.
Table~\ref{tab:sh_ranges} summarises the search ranges considered for each
optimisation method.
Configuration selection is performed using a two-stage successive halving (SH)
procedure, in which candidate configurations are progressively evaluated under
increasing computational budgets and poorer configurations are discarded.

\begin{table}[t]
\centering
\caption{Hyperparameter search ranges used in Stage~II for each optimiser.}
\label{tab:sh_ranges}
\begin{adjustbox}{max width=\linewidth}
\begin{tabular}{ll}
\toprule
\textbf{Optimiser} & \textbf{Hyperparameter Search Range} \\
\midrule
Adam &
Learning rate $\in \{10^{-4}, 3\times10^{-4}, 10^{-3}, 3\times10^{-3}, 3\times10^{-2}\}$;
AMSGrad $\in \{\text{True}, \text{False}\}$ \\

RMSprop &
Learning rate $\in \{10^{-4}, 3\times10^{-4}, 10^{-3}\}$;
$\alpha \in \{0.9, 0.99\}$ \\

SGD &
Learning rate $\in \{10^{-4}, 3\times10^{-4}, 10^{-3}, 3\times10^{-3}\}$;
Momentum $\in \{0.0, 0.9\}$ \\

PSO &
Inertia weight $w \in \{0.4, 0.72, 0.9\}$;
$c_1, c_2 \in \{1.2, 1.49, 2.0\}$ \\

GA &
Mutation rate $\in \{0.05, 0.1, 0.2\}$ \\

DE &
Scaling factor $F \in \{0.5, 0.8, 1.0\}$;
Crossover rate $CR \in \{0.7, 0.9\}$ \\

CMA-ES &
Initial step size $\sigma \in \{0.1, 0.3, 0.5\}$ \\
\bottomrule
\end{tabular}
\end{adjustbox}
\end{table}

Table~\ref{tab:sh_selected} reports the optimiser-specific configurations
selected at the final stage of the successive halving procedure.
These configurations are fixed for all subsequent experiments.

\begin{table}[t]
\centering
\caption{Optimiser configurations selected by successive halving (Stage~II).}
\label{tab:sh_selected}
\begin{adjustbox}{max width=\linewidth}
\begin{tabular}{ll}
\toprule
\textbf{Optimiser} & \textbf{Selected Configuration} \\
\midrule
Adam &
Learning rate = 0.03; AMSGrad = False; Weight decay = 0.0 \\

SGD &
Learning rate = 0.003; Momentum = 0.9; Nesterov = False \\

RMSprop &
Learning rate = 0.001; $\alpha = 0.99$; Momentum = 0.0 \\

PSO &
Population size = 30; Iterations = 10;
$w = 0.9$; $c_1 = 1.2$; $c_2 = 1.2$ \\

GA &
Population size = 30; Iterations = 10;
Mutation rate = 0.05 \\

DE &
Population size = 30; Iterations = 10;
$F = 0.5$; $CR = 0.9$ \\

CMA-ES &
Population size = 30; Iterations = 10;
$\sigma = 0.1$ \\
\bottomrule
\end{tabular}
\end{adjustbox}
\end{table}

% ===================================================
\subsection{Stage III: Final Evaluation}

Final evaluation is conducted using the optimiser-specific configurations
selected in Stage~II.
Each optimiser is evaluated on all ten datasets using ten fixed random seeds
under identical experimental conditions.
The test set is accessed exclusively at this stage and only once per
configuration.

Table~\ref{tab:stage3_results} summarises the final test-set performance obtained
by each optimiser.
Due to space constraints, only the table structure is shown here; numerical
results are reported and discussed in the following section.

\begin{table*}[t]
\centering
\caption{Final test RMSE obtained by each optimiser across ten datasets
(best over ten seeds).}
\label{tab:stage3_results}
\begin{adjustbox}{max width=\textwidth}
\begin{tabular}{lccccccc}
\toprule
Dataset & Adam & RMSprop & SGD & PSO & GA & DE & CMA-ES \\
\midrule
Abalone   
& \textbf{2.118} & 2.221 & 2.228 & 2.373 & 2.788 & 2.703 & 3.414 \\

Airfoil  
& 6.987 & 7.797 & 7.444 & 7.408 & 7.948 & \textbf{6.863} & 7.689 \\

Ankara   
& \textbf{2.405} & 6.321 & 4.627 & 2.805 & 3.920 & 3.253 & 5.380 \\

AutoMPG6 
& \textbf{2.267} & 2.379 & 2.386 & 2.628 & 2.651 & 2.544 & 2.673 \\

Baseball 
& \textbf{726.507} & 793.069 & 774.236 & 776.944 & 760.099 & 758.258 & 802.368 \\

Concrete 
& \textbf{10.877} & 12.128 & 11.672 & 11.262 & 11.368 & 11.423 & 12.208 \\

Izmir    
& \textbf{2.052} & 4.154 & 2.811 & 2.319 & 3.308 & 2.853 & 3.776 \\

Laser    
& \textbf{5.445} & 6.219 & 5.633 & 6.694 & 8.319 & 7.827 & 9.462 \\

Treasury 
& 1.195 & 1.223 & \textbf{0.989} & 1.187 & 1.148 & 1.156 & 1.283 \\

Wine     
& \textbf{0.097} & 0.611 & 0.592 & 0.206 & 0.338 & 0.211 & 0.486 \\
\bottomrule
\end{tabular}
\end{adjustbox}
\end{table*}




% ===================================================
% ===================================================
\section{Discussion}

The Stage~III results provide a controlled comparison of gradient-based and
non-gradient optimisation methods for Mamdani-type fuzzy inference systems under
a unified experimental protocol.
Since all optimisers are evaluated using configurations selected in Stage~II,
with identical model structure, initialisation strategy, and optimisation
budgets, the observed performance differences can be attributed primarily to
intrinsic optimiser behaviour rather than confounding experimental factors.

A first key observation is that no single optimisation method dominates across
all datasets.
Nevertheless, clear trends emerge when examining performance across the ten
benchmark problems.
Among the methods considered, Adam achieves the lowest test RMSE on the largest
number of datasets, including Abalone, Ankara, AutoMPG6, Baseball, Concrete,
Izmir, Laser, and Wine (Table~\ref{tab:stage3_results}).
This indicates that adaptive gradient-based optimisation provides a strong and
robust default choice for a wide range of regression tasks under the evaluated
experimental conditions.

Non-gradient and population-based optimisers remain competitive on specific
datasets.
Differential Evolution (DE) achieves the best performance on Airfoil, while SGD
outperforms all other methods on Treasury.
Particle Swarm Optimisation (PSO) consistently attains near-optimal performance
across several datasets, often ranking second-best, although it does not achieve
the lowest RMSE as frequently as Adam.
These results suggest that population-based methods can offer advantages on
particular problem instances, but their benefits are not uniform across
datasets.

Another notable finding is the variability in optimiser behaviour across
datasets with differing characteristics.
Datasets for which gradient-based methods perform strongly tend to exhibit
smoother optimisation landscapes, where local gradient information can be
exploited effectively.
In contrast, the relative success of population-based methods on some datasets
suggests increased robustness to irregular or multimodal error surfaces, albeit
at the cost of greater computational effort.

Overall, the observed performance differences highlight that optimiser behaviour
in fuzzy inference systems is strongly dataset-dependent.
While adaptive gradient-based methods benefit from smooth or moderately
structured loss landscapes, population-based approaches exhibit greater
robustness to irregular or multimodal surfaces.
These findings underscore the importance of controlled experimental protocols
for disentangling optimiser behaviour from structural and configuration-related
effects.


\section{Conclusion}

This paper proposed a staged and reproducible experimental framework for the
controlled comparison of gradient-based and non-gradient optimisation methods
for Mamdani-type fuzzy inference systems.
By explicitly decoupling structural design decisions, optimiser configuration
selection, and final evaluation, the framework reduces methodological
confounding and enables interpretable, system-level analysis of optimiser
behaviour.

Empirical results across ten benchmark regression datasets show that no single
optimisation method is universally optimal.
However, adaptive gradient-based optimisation, particularly Adam, demonstrates
the strongest overall robustness under the fixed experimental protocol,
achieving the best performance on the largest number of datasets when combined
with the configurations selected via successive halving.

Future work will extend the proposed framework to larger benchmark suites,
broader configuration spaces, and alternative fuzzy inference architectures,
enabling more comprehensive evaluation of optimisation strategies under
increased experimental coverage.



% ===================== REFERENCES =====================
\begin{thebibliography}{99}

\bibitem{zadeh1965}
L.~A. Zadeh, ``Fuzzy sets,'' \emph{Information and Control}, vol.~8, no.~3,
pp.~338--353, 1965.

\bibitem{holland1975ga}
J.~H. Holland, \emph{Adaptation in Natural and Artificial Systems}.
Ann Arbor, MI, USA: University of Michigan Press, 1975.

\bibitem{kennedy1995pso}
J.~Kennedy and R.~Eberhart, ``Particle swarm optimization,'' in
\emph{Proc. IEEE Int. Conf. Neural Networks (ICNN)}, vol.~4, 1995,
pp.~1942--1948.

\bibitem{storn1997de}
R.~Storn and K.~Price, ``Differential evolution -- a simple and efficient
heuristic for global optimization over continuous spaces,''
\emph{Journal of Global Optimization}, vol.~11, no.~4, pp.~341--359, 1997.

\bibitem{hansen2001cmaes}
N.~Hansen and A.~Ostermeier, ``Completely derandomized self-adaptation in
evolution strategies,'' \emph{Evolutionary Computation}, vol.~9, no.~2,
pp.~159--195, 2001.

\bibitem{kingma2015adam}
D.~P. Kingma and J.~Ba, ``Adam: A method for stochastic optimization,'' in
\emph{Proc. Int. Conf. Learning Representations (ICLR)}, 2015.

\bibitem{tieleman2012rmsprop}
T.~Tieleman and G.~Hinton, ``Lecture 6.5---RMSprop: Divide the gradient by a
running average of its recent magnitude,'' \emph{Neural Networks for Machine
Learning}, Coursera, 2012.

\bibitem{jamieson2016sh}
K.~Jamieson and A.~Talwalkar, ``Non-stochastic best arm identification and
hyperparameter optimization,'' in
\emph{Proc. 19th Int. Conf. Artificial Intelligence and Statistics (AISTATS)},
2016.

\bibitem{chen2024gradient}
C.~Chen, C.~Wagner, and J.~M. Garibaldi,
``Gradient-based fuzzy system optimisation via automatic differentiation --
FuzzyR as a use case,'' 2024.

\bibitem{zhou2023stochastic}
W.~Zhou, H.~Li, and M.~Bao,
``Stochastic configuration based fuzzy inference system with interpretable fuzzy
rules and intelligent search process,''
\emph{Mathematics}, vol.~11, no.~3, 2023.

\bibitem{chen2025casp}
C.~Chen, J.~M. Mendel, and J.~M. Garibaldi,
``Constraints always satisfied parameters (CASPs) for fuzzy sets optimisation,''
\emph{IEEE Transactions on Fuzzy Systems}, in press, 2025.

\bibitem{cordon2001review}
O.~Cordón, F.~Herrera, F.~Hoffmann, and L.~Magdalena,
``A review on the combination of evolutionary algorithms and fuzzy systems,''
\emph{Information Sciences}, vol.~136, no.~1--4, pp.~1--38, 2001.


\end{thebibliography}

\end{document}

